\newcommand{\DoNotLoadEpstopdf}{}
\documentclass[11pt,reqno]{amsart}
\usepackage[T1]{fontenc}
\usepackage{lmodern,amsmath,amssymb,amsthm,mathtools,mathrsfs,microtype}
\usepackage[margin=1.1in]{geometry}
\usepackage{xr-hyper}
\usepackage[hidelinks]{hyperref}
\externaldocument[supp-]{build/supplement}[supplement.pdf]
\hypersetup{pdftitle={Collision records, arithmetic, and raw local inversion in a triangular Lorentz gas},pdfauthor={Qian Qi}}
\setlength{\emergencystretch}{2em}
\allowdisplaybreaks[2]
\numberwithin{equation}{section}
\newtheorem{theorem}{Theorem}[section]
\newtheorem{lemma}[theorem]{Lemma}
\newtheorem{proposition}[theorem]{Proposition}
\newtheorem{corollary}[theorem]{Corollary}
\theoremstyle{definition}
\newtheorem{definition}[theorem]{Definition}
\newtheorem{inputstatement}[theorem]{Input statement}
\theoremstyle{remark}\newtheorem{remark}[theorem]{Remark}
\newcommand{\R}{\mathbb R}\newcommand{\Z}{\mathbb Z}
\newcommand{\one}{\mathbf1}\newcommand{\dd}{\,\mathrm d}
\newcommand{\esssup}{\operatorname*{ess\,sup}}
\newcommand{\diver}{\operatorname{div}}
\newcommand{\TV}{\operatorname{TV}}
\title[Collision records and raw local inversion]{Collision records, arithmetic, and raw local inversion in a triangular Lorentz gas}
\author{Qian Qi}
\date{October 12, 2026. A2-DYN, revision 90}
\subjclass[2020]{37D50, 37A50, 60F05, 28A75}
\keywords{Lorentz gas, exact collision records, microscopic oscillation, original-source regularity, critical entrances, exact local limits}
\begin{document}
\raggedbottom
\begin{abstract}
We retain the pointwise local limit problem for the complete displacement,
collision-count, occupation and roof record of actual returns in a
triangular Lorentz gas. We prove finite-count oscillation estimates on
the original collision source, including both true physical ends and
backward critical entrances. A capped end-layer estimate and radial
motion yield a uniform $L^1$ translation modulus for the full weighted
roof density. A one-endpoint-per-segment argument then gives a modulus
of order $h^{1/33}$ for every measurable initial segment restriction,
uniformly in the endpoint location. This includes the original union
of initial criticality and short total residence, with both cutoffs
arbitrarily small. Its complete microscopic fractional-variation tail
is bounded before the collision limit, and its fractional local norm
has a quantitative vanishing bound. All exact labels, physical weights,
arithmetic classes and previous proofs are retained. The new estimates
control integrated roof oscillation; the unchanged essential-height
charge still requires a uniform starting-roof estimate.
\end{abstract}
\maketitle

\input{journal-problem-v89}

\subsection{An all-scale estimate on the original exceptional source}
Let $\tau^-,\tau^+$ be the stopped roof-flow lifetimes, $\ell=\tau^-+\tau^+$
and $z=|\nabla_qF|$. The original exceptional source is
\[
 R_{a,\delta}^\varepsilon(t)
   =\int W^\varepsilon(x)\one_{\{z(x)<a\ \mathrm{or}\ \ell(x)<\delta\}}
                                                       \,d\mu_t(x).
\]
Its union is counted once. Neither endpoint section membership nor a
true boundary is removed. Our principal new assertion is the following.
\begin{theorem}[Microscopic oscillation on the complete original union]
\label{thm:v90-main}
Under the precise physical and spectral inputs stated below, for every
fixed $L$ and $0<\varepsilon\le1$ there is $K_L<\infty$ such that
\[
 \sup_{m\ge2,R,n,k}\sup_{a,\delta>0}\sup_{|J|\le L}m^2
 \int_J|R_{a,\delta}^\varepsilon(t+h)-R_{a,\delta}^\varepsilon(t)|\,dt
       \le K_L\bigl(h^{1/33}+\varepsilon^{-1}\sqrt h\bigr)
\]
for every sufficiently small $h>0$. The constant is obtained at one
fixed auxiliary band. The supremum over finite collision counts and
both cutoffs precedes $h\downarrow0$.
For $0<s<1/33$, the full integral of these differences against
$h^{-1-s}dh$ down to zero has a uniform vanishing microscopic tail.
For the cofinal choice $a=r^{1/8}$, $\delta=r^{1/4}$,
\[
 \limsup_m\sup_{R,n,k}\sup_{|J|=L}
 \|m^2R_{a,\delta}^\varepsilon\|_{s,1;J}
             \le C_{L,s}(1+\varepsilon^{-1})^{33s}r^{(1-33s)/128},
\]
where the fractional norm is defined in \eqref{eq:v90-gagliardo}.
\end{theorem}
The proof supplies a genuine pre-limit envelope, rather than a sum of
fixed-scale limiting estimates. The complete density $b^\varepsilon$
has the stronger translation exponent $1/32$. The loss to $1/33$ for
the original union comes from the single new interior cutoff on a
physical segment. Its location needs no regularity in angle.

These are oscillation statements in the roof variable, not another
claim that a local mass estimate is a pointwise bound. In particular the
integral over starting roofs is visible in the theorem. The exact
pointwise objective remains \eqref{eq:v76-target}; no unconditional
point-conditioned law is deduced without its essential-height closure.

\input{journal-inputs-v90}
\input{core/211_two_sided_end_layers_and_source_oscillation}
\input{core/212_cutoff_uniform_original_source_prefixes}
\input{core/213_microscopic_fractional_variation_and_endpoint_scope}

\section{Dependency and scalar-first order}
The new argument has three steps. The two-sided end-layer theorem pays
all true boundary changes, including critical backward ends. The
weighted segment calculation controls surviving variation without an
initial inverse-gradient cutoff. The prefix theorem pays the extra
cut introduced by the original exceptional union without a derivative
of its angular graph. Every estimate is summed over actual physical
interiors using the same Borel disintegration.

The preceding critical-distribution and diffusive-window theorems remain
in the complete preceding article in the supplement. They are not
replaced, and the two revision-88 variants remain separately preserved.
The new physical oscillation envelope does not establish zero
essential-height charge. A further uniform starting-roof or comparable
pointwise persistence bound is still required for that conclusion.
Only after scalar closure may the retained positive-arithmetic
normalization, same-roof bridge, clock transfer and likelihood arguments
be used for unconditional point-conditioned laws.

\input{journal-references-v81}
\end{document}
